\subsection{An explicit new unit and its three independent conjugate classes}

\begin{statement}{A.197}{Proved}
{A literal ninth-root-field unit reduces the subfield-unit quotient to order at most 27}
Retain the literal fields, roots, units, and subgroup \(B\) of
A.191 and A.194.  Thus
\[
 \alpha^3=3,\quad d=649,\quad t^3=d^2\alpha,\quad
 z=\zeta_9,\quad \zeta=z^3,\quad U=\alpha^2-2,
 \qquad M=\Q(t,z).
\tag{A.197.1}
\]
Use exactly the units \(e_1,\ldots,e_9\) of A.191 and
\(u_2,u_4\) of A.159.  The element
\(\alpha u_2u_4/U\) has a unique real cube root \(r_0\)
in the actual embedded field \(K_2=\Q(t)\).
Define
\[
\begin{aligned}
 c={}&\frac{-1+z^3}{3}
       +\frac{\alpha(1+z+z^4)}{3}\\
    &-\frac{\alpha^2(2+2z+2z^2+z^3+z^4+z^5)}{9},
       \qquad r=cr_0.
\end{aligned}
\tag{A.197.2}
\]
Then \(r\) is an actual global unit of \(M\), with the
literal identity
\[
                r^3e_3e_6e_9=e_1e_5e_8u_2u_4.
\tag{A.197.3}
\]
Let \(g(t)=zt\), \(g(z)=z\), as in A.194.  In the
pro-three completion \(\mathcal U_M\), let \(B^+\) be
the closed subgroup generated by \(B\) and the nine actual
units \(g^j(r)\), \(0\leq j<9\).  Then
\[
       B^+/B\simeq(\Z/3)^3,\qquad
       \mathcal U_M/B^+\simeq(\Z/3)^b,
                    \quad 0\leq b\leq3.
\tag{A.197.4}
\]
The three classes of \(r,g(r),g^2(r)\) form the displayed
basis of \(B^+/B\).  In particular
\(3\leq\dim_{\F_3}(\mathcal U_M/B)\leq6\).
The image of \(B^+\) in \(\mathcal U_M/3\mathcal U_M\)
has dimension at least 24.  The displayed bound uses an
explicit nonzero minor and leaves its full dimension unspecified.
\end{statement}

\begin{proof}
We first prove the field inclusion used in the definition of
\(r_0\).  Put \(\gamma_{11}=2+\alpha\).  The complete
S-unit detector of A.165 gives the following two identities
in the actual quotient \(K_2^\times/(K_2^\times)^3\):
\[
 \frac{u_2}{U^{19}}\equiv\gamma_{11}^2\,11^2\,59,
 \qquad
 \frac{u_4}{U^{538}}\equiv U^2\gamma_{11}\,59.
\tag{A.197.5}
\]
For precision, order the base generators as
\((U,\alpha,\alpha-1,\gamma_{11},
2+2\alpha-\alpha^2,2,11,59)\).
The two right sides have exponent vectors
\((0,0,0,2,0,0,2,1)\) and
\((2,0,0,1,0,0,0,1)\), respectively.
Evaluation in each of the twelve literal residue maps of
A.165 gives the same vector as the corresponding left side.
A.165 proves that those maps are injective on the entire
S-unit quotient modulo cubes.  Thus this comparison proves
(A.197.5) in that quotient, with its actual restriction map
from the base field.  The certificate
\path{code/two_inert_base_restriction_unit_classes.py}
retains all twelve comparisons and their rank-twelve test.

Multiplying (A.197.5) and using \(d=11\cdot59\) gives
\[
 \frac{\alpha u_2u_4}{U}
       \equiv\alpha d^2U^{558}\gamma_{11}^3
       =(tU^{186}\gamma_{11})^3
            \quad\bmod (K_2^\times)^3.
\tag{A.197.6}
\]
Therefore the left side has a cube root in \(K_2\).
This field has odd degree and has its specified real
embedding, so it contains no nontrivial cube root of unity.
The root is consequently unique in \(K_2\), and equals
the unique real cube root stated above.

Using only \(\alpha^3=3\), \(z^6+z^3+1=0\), and the
nine literal A.191 definitions, multiplication of the small
expression (A.197.2) gives
\[
          c^3\alpha e_3e_6e_9=Ue_1e_5e_8.
\tag{A.197.7}
\]
Equations (A.197.6)--(A.197.7) prove (A.197.3) as an
equality of elements of \(M\).  Its right side divided
by \(e_3e_6e_9\) is a product of certified global units.
Thus both \(r\) and \(r^{-1}\) satisfy monic cubic
equations with algebraic-integer coefficients.  Both are
algebraic integers, so \(r\in\mathcal O_M^\times\).

There is a separate rational-coordinate verification of
these identities.  The retained 54 coefficients of a root
\(\widetilde r\) are in the basis
\(t^iz^j\), \(0\leq i<9\), \(0\leq j<6\).
Exact rational multiplication verifies (A.197.3) for that
root.  Inverting the small element \(c\) in
\(F'=\Q(\alpha,z)\), the same calculation gives
\(\zeta^2\widetilde r/c\in\Q(t)\) and verifies its
cube to be \(\alpha u_2u_4/U\).  Uniqueness therefore
identifies this element with \(r_0\), and proves
\(\widetilde r=\zeta r\).  Since \(\zeta=z^3\),
the two literal roots have exactly the same class modulo
cubes in \(M\).  The same equality of cube classes holds
after every \(g^j\).  This comparison makes the compact
definition and the independent rational certificate agree.

For the rank calculation, take the 33 units in the precise
order (A.194.6), followed by
\[
                       r,g(r),\ldots,g^8(r).
\tag{A.197.8}
\]
At each row \((p,a,b)\) below, use the actual residue map
\((t,z)\mapsto(a,b)\).  Define an entry \(m_i\in\{0,1,2\}\)
by \(w_i(a,b)^{(p-1)/3}=b^{3m_i}\), where \(w_i\)
is the corresponding unit in (A.197.8).
All denominators are invertible and all unit values are
nonzero in the displayed residue fields.
\[
\begin{array}{r|r|r@{\qquad}r|r|r@{\qquad}r|r|r}
 p&a&b&p&a&b&p&a&b\\\hline
307&78&287&307&78&93&307&78&53\\
307&78&168&307&78&274&307&78&46\\
307&98&287&307&98&93&307&98&53\\
307&131&287&307&143&287&307&143&93\\
307&143&53&307&189&287&523&20&468\\
523&20&410&523&20&217&523&154&468\\
523&154&410&523&156&468&523&156&410\\
523&311&468&523&311&410&613&47&467
\end{array}
\tag{A.197.9}
\]
Read the rows across each line.  In this 24-by-42 integer
matrix, the columns
\[
\begin{split}
 1,2,3,4,5,6,7,8,9,11,12,15,16,17,19,20,23,24,25,\\
                         31,32,34,35,36
\end{split}
\tag{A.197.10}
\]
have the exact determinant
\[
                       181379512\equiv1\pmod3.
\tag{A.197.11}
\]
The literal formula for \(g\) evaluates a conjugate by
substituting \(t\mapsto tb^j\), \(z\mapsto b\).
Thus this determinant uses the actual conjugate units and
gives a rank of at least 24 in the completed unit quotient.
On these same 24 rows, all 33 old columns together have
rank exactly 21, while the old columns together with
\(r,g(r),g^2(r)\) have rank 24.  Consequently no nonzero
\(\F_3\)-linear combination of those three classes belongs
to \(B\): reduction would put the same combination in
the span of the entire old column set.

The unique prime \(P_3\) of \(K_2\) above three satisfies
\(\operatorname{div}_{K_2}(\alpha)=3P_3\).  Since
\(r_0^3=\alpha u_2u_4/U\), the exact divisor map gives
\(\operatorname{div}_{K_2}(r_0)=P_3\).
In \(E_2=K_2(\zeta)\), its extension is fixed by
\(\sigma=g^3\).  Hence
\[
             \frac{g^3(r)}{r}
                   =\frac{\sigma(r_0)}{r_0}
                       \in\mathcal O_{E_2}^{\times}.
\tag{A.197.12}
\]
Here \(g^3\) fixes \(c\in F'\), so the displayed
comparison is literal.  The automorphism \(g\) preserves
\(F'\) and cyclically permutes the three fields \(E_{2,j}\).
It therefore preserves \(B\).  Applying its powers to
(A.197.12) proves that all nine orbit classes are generated
modulo \(B\) by the first three.  Their independence just
proved, and the exponent-three assertion of A.194, give
the first isomorphism in (A.197.4).

A.194 proves \(3\mathcal U_M\subseteq B\subseteq B^+\)
and \(\dim_{\F_3}\mathcal U_M/3\mathcal U_M=27\).
The exact quotient sequence
\[
 0\longrightarrow B^+/3\mathcal U_M
 \longrightarrow\mathcal U_M/3\mathcal U_M
 \longrightarrow\mathcal U_M/B^+\longrightarrow0
\tag{A.197.13}
\]
therefore proves (A.197.4).

The fixed certificate
\path{code/two_inert_ninth_root_orbit_unit_certificate.py}
checks the rational 54-coordinate identity, the compact
18-coordinate identity, their exact field comparison, all
1008 residue values, determinant (A.197.11), and both full
column-set ranks used above.  Its output
and data are retained.  The earlier bounded discovery search
is retained separately and supplies no upper bound on the
dimension of the image of the generated units.
\end{proof}
